\chapter{Appendices}\label{sec:appendices}

\section{Appendix A: Mathematical Notation}\label{sec:appendix-a-mathematical-notation}

{\def\LTcaptype{table} % do not increment counter
\begin{longtable}[]{@{}ll@{}}
\toprule\noalign{}
Symbol & Description \\
\midrule\noalign{}
\endhead
\bottomrule\noalign{}
\endlastfoot
\(b\) & Binary bit (\(\in \{0, 1\}\)) \\
\(x\) & Modulated BPSK symbol (\(\in \{-1, +1\}\)) \\
\(y\) & Received signal \\
\(n\) & Noise sample \\
\(h\) & Fading coefficient \\
\(\sigma^2\) & Noise variance \\
\(\text{SNR}\) & Signal-to-Noise Ratio (linear) \\
\(G\) & AF amplification gain \\
\(\mathbf{H}\) & MIMO channel matrix (\(\in \mathbb{C}^{2 \times 2}\)) \\
\(\mathbf{W}\) & Neural network weight matrix \\
\(\mathbf{b}\) & Bias vector \\
\(P_e\) & Bit error rate \\
\(Q(\cdot)\) & Gaussian Q-function \\
\(K\) & Rician K-factor \\
\(\beta\) & VAE KL weight \\
\(\lambda\) & Regularization parameter \\
\(\eta\) & Learning rate \\
\(\Delta\) & State space discretization step \\
\(\mathbf{A}, \mathbf{B}, \mathbf{C}, D\) & State space model matrices \\
\end{longtable}
}

\section{Appendix B: Model Architectures and Hyperparameters}\label{sec:appendix-b-model-architectures-and-hyperparameters}

\textbf{MLP (Minimal):} - Input: 5-symbol sliding window - Hidden: 24 neurons, ReLU activation - Output: 1 neuron, Tanh activation - Parameters: 169 - Training: MSE loss, lr=0.01, 100 epochs, 25K samples at SNR={[}5, 10, 15{]} dB - Implementation: NumPy (CPU)

\textbf{Hybrid:} - Architecture: Same as MLP (169 params) - SNR threshold: Learned (default \textasciitilde5 dB) - Below threshold → MLP processing; Above threshold → DF processing - Implementation: NumPy (CPU)

\textbf{VAE:} - Encoder: 7 → 32 (ReLU) → 16 (ReLU) → μ(8), log σ²(8) - Decoder: 8 → 16 (ReLU) → 32 (ReLU) → 1 (Tanh) - Parameters: 1,777 - Training: β-VAE loss (β=0.1), Adam lr=1e-3, 100 epochs - Implementation: PyTorch (CUDA)

\textbf{CGAN (WGAN-GP):} - Generator: (7+8) → 32 (LeakyReLU) → 32 (LeakyReLU) → 16 (LeakyReLU) → 1 (Tanh) - Critic: (1+7) → 32 (LeakyReLU) → 16 (LeakyReLU) → 1 - Parameters: 2,946 - Training: WGAN-GP (λ\_GP=10, λ\_L1=100), Adam lr=1e-4, 200 epochs, 5 critic updates per generator update - Implementation: PyTorch (CUDA)

\textbf{Transformer:} - Input projection: 1 → 32 - Positional encoding: Sinusoidal - Encoder: 2 layers, 4 heads, d\_model=32, d\_ff=128 - Output projection: 32 → 1 (Tanh) - Parameters: 17,697 - Training: MSE loss, Adam lr=1e-3, 100 epochs - Implementation: PyTorch (reported experiments on CUDA; CPU fallback supported)

\textbf{Mamba S6:} - Input projection: 1 → 32 - Mamba blocks: 2 layers, each: LayerNorm → expand (32→64) → S6 (d\_state=16) → contract (64→32) → residual - S6 selective parameters: Δ, B, C = f(input) via learned linear projections - Output projection: 32 → 1 (Tanh) - Parameters: 24,001 - Training: MSE loss, Adam lr=1e-3, 100 epochs - Implementation: PyTorch (reported experiments on CUDA; CPU fallback supported)

\textbf{Mamba2 (SSD):} - Input projection: 1 → 32 - Mamba-2 blocks: 2 layers, each: LayerNorm → parallel branches (SiLU gate ∥ SSD layer) → gated output → contract (64→32) → residual - SSD layer: chunk\_size=8, builds lower-triangular causal kernel \(M \in \mathbb{R}^{L \times L}\) per chunk via cumulative log-decay, applies \(Y = M \cdot V\) as batched matmul - Inter-chunk state: running state \((B, N, D)\) updated once per chunk - S4D-style A initialisation: \(A_{\log} = \log(1, 2, \dots, d_{\text{state}})\) - Selective parameters: Δ, B, C = f(input); value projection V = f(input) - Output projection: 32 → 16 (SiLU) → 1 (Tanh) - Parameters: 26,179 - Training: MSE loss, Adam lr=1e-3, 100 epochs, gradient clipping (max\_norm=1.0) - Implementation: PyTorch (CUDA)

\section{Appendix C: Software Architecture}\label{sec:appendix-c-software-architecture}

The project is implemented as a modular Python package (\texttt{relaynet}) with the following structure:

\begin{verbatim}
relaynet/
├── channels/          # Channel models
│   ├── awgn.py            # AWGN channel
│   ├── fading.py          # Rayleigh & Rician fading
│   └── mimo.py            # 2×2 MIMO + ZF/MMSE/SIC equalization
├── modulation/
│   ├── bpsk.py            # BPSK modulate/demodulate
│   ├── qpsk.py            # QPSK Gray-coded modulate/demodulate
│   ├── qam.py             # 16-QAM Gray-coded modulate/demodulate
│   └── psk.py             # 16-PSK modulate/demodulate
├── relays/            # Relay strategies
│   ├── base.py            # Abstract base class
│   ├── af.py              # Amplify-and-Forward
│   ├── df.py              # Decode-and-Forward
│   ├── genai.py           # Minimal MLP (feedforward NN)
│   ├── hybrid.py          # SNR-adaptive Hybrid
│   ├── vae.py             # Variational Autoencoder
│   ├── cgan.py            # Conditional GAN (WGAN-GP)
│   ├── transformer.py     # Transformer relay
│   ├── mamba.py           # Mamba S6 relay
│   └── mamba2.py          # Mamba-2 SSD relay

├── simulation/
│   ├── runner.py          # Monte Carlo BER simulation
│   └── statistics.py      # CI computation, significance tests
├── visualization/
│   └── plots.py           # BER plotting utilities
└── utils/
    └── torch_compat.py    # Device detection helpers
\end{verbatim}

The framework uses object-oriented design with a common \texttt{Relay} base class, enabling polymorphic relay swapping. Monte Carlo simulation is implemented in \texttt{runner.py} with configurable trial count, bit count, and SNR range. All MIMO operations use vectorized PyTorch for GPU acceleration.

\textbf{Testing:} 126 automated tests (pytest) cover all channels, modulation (BPSK, QPSK, 16-QAM, 16-PSK), relay strategies, simulation, statistics, and modulation-comparison modules with 100\% pass rate.

\textbf{Reproducibility:} Random seeds are controlled at the source (bit generation) and noise (per-trial seeding) levels to ensure reproducible results.

\section{Appendix D: Normalized 3K-Parameter Configurations}\label{sec:appendix-d-normalized-3k-parameter-configurations}

{\def\LTcaptype{table} % do not increment counter
\begin{longtable}[]{@{}
  >{\raggedright\arraybackslash}p{(\linewidth - 6\tabcolsep) * \real{0.2500}}
  >{\raggedright\arraybackslash}p{(\linewidth - 6\tabcolsep) * \real{0.2500}}
  >{\raggedright\arraybackslash}p{(\linewidth - 6\tabcolsep) * \real{0.2500}}
  >{\raggedright\arraybackslash}p{(\linewidth - 6\tabcolsep) * \real{0.2500}}@{}}
\toprule\noalign{}
\begin{minipage}[b]{\linewidth}\raggedright
Model
\end{minipage} & \begin{minipage}[b]{\linewidth}\raggedright
Parameters
\end{minipage} & \begin{minipage}[b]{\linewidth}\raggedright
Window
\end{minipage} & \begin{minipage}[b]{\linewidth}\raggedright
Hidden / Architecture
\end{minipage} \\
\midrule\noalign{}
\endhead
\bottomrule\noalign{}
\endlastfoot
MLP-3K & 3,004 & 11 & hidden=231 \\
Hybrid-3K & 3,004 & 11 & hidden=231 (+ DF switch) \\
VAE-3K & 3,037 & 11 & latent=10, hidden=(44, 20) \\
CGAN-3K & 3,004 & 11 & noise=8, g\_hidden=(30, 30, 16), c\_hidden=(32, 16) \\
Transformer-3K & 3,007 & 11 & d\_model=18, heads=2, layers=1 \\
Mamba-3K & 3,027 & 11 & d\_model=16, d\_state=6, layers=1 \\
Mamba2-3K & 3,004 & 11 & d\_model=15, d\_state=6, chunk\_size=8, layers=1 \\
\end{longtable}
}

All 3K configurations use a window size of 11 (vs.~5 for original MLP/Hybrid, and 11 for original sequence models) to provide a common input context. The parameter counts are within ±1.2\% of the 3,000 target.

\begin{center}\rule{0.5\linewidth}{0.5pt}\end{center}

%% ── Hebrew Abstract (Back Matter) ───────────────────────────
